\section{Discussion}

The choice of lattice basis affects a drifting model at the point where it compares a generated crystal with its references. Using physical distances and angles in a periodic neighborhood makes this comparison independent of the coordinate convention used for each crystal. The same field can supervise a network that moves individual atoms or a network that moves whole Wyckoff orbits. These generators make different predictions even when they use the same definition of distance. In particular, an orbit update preserves the operations of the selected group, so the choice of occupied pattern limits which structures it can reach.

The practical advantage of one-step sampling depends on the time available and on the starting structure. With random starting sites, the compact sixteen-step flow recovered slightly more targets per twenty attempts. Under a short time limit, the site generator could try more starting configurations and recovered more targets than the flow. Returning the sampled starting structures was faster still and yielded the most matches on that panel. A stronger prior changes the comparison. On the same symbolic seeds, a learned orbit update improved five-draw recovery by just over four percentage points. But when the catalog could supply an entire training geometry, transferring that geometry was often a better prediction than drawing new coordinates inside the retrieved orbits. On the protected family holdout the catalog contained none of the reference occupation patterns, and the improvement over the seeds was small. Whether a structural family is represented in training matters when interpreting recovery from a formula.

The matched training runs also help locate the source of the improvement. Contact features and angular features gave almost identical recovery rates, and a MeanFlow objective trained on the same backbone and seeds came close to both. The clearest evidence for the periodic drift is therefore that it provides an invariant training signal and learns a useful update. The measured speed advantage over the compact iterative flow belongs to the one-step sampling procedure. The comparison with the released one-step uFlowCSP model illustrates the distinction: the orbit generator produced proposals faster on the same hardware, while uFlowCSP recovered more targets at a fixed proposal count.

Two limitations matter for materials discovery. First, a structural match that allows the lattice to rescale can conceal an error in the predicted density. Restoring the seed volume while retaining the learned shape and sites improved volume-qualified recovery, which points to scale prediction as a distinct source of error. Second, sampling compositions from the training distribution tests new arrangements at familiar compositions. It gives no estimate of the probability of discovering a stable material at a previously unseen composition. For a search with a fixed compute budget, the relevant outcome is how many candidates survive structural matching, relaxation, and stability and novelty checks within that budget. Here the measured speedup concerns the proposal stage.

\paragraph{Data and code.}
The experiments use the MP-20 partition of Materials Project structures~\cite{materialsproject,cdvae}. The project records the training and evaluation splits, the ordered proposals, the selected checkpoints, and the code for computing the reported metrics. External model weights are identified with their respective methods. The training and evaluation code is available at \url{https://github.com/AMIRMOHAMMAD-OSS/CrystalDrift}.

\paragraph{Acknowledgment.}
OpenAI Codex assisted with software, experiments, verification, and manuscript preparation.
